\introChapter{เมทริกซ์และดีเทอร์มิแนนต์ (Matrices \& Determinants)}

\begin{description}[leftmargin=2cm, style=nextline, itemsep=0.1em, parsep=0pt]
    \item[รหัสวิชา] 4091606
    \item[ชื่อวิชา] คณิตศาสตร์สำหรับคอมพิวเตอร์
    \item[หน่วยกิต] 3 (3-0-6)
    \item[บทที่] 5
    \item[ชื่อบท] เมทริกซ์และดีเทอร์มิแนนต์ (Matrices \& Determinants)
    \item[จำนวนชั่วโมง] 9 ชั่วโมง
\end{description}

\noindent \textbf{\large วัตถุประสงค์การเรียนรู้ประจำบท}\par\vspace{0.2em}\noindent
\noindent เมื่อนักศึกษาเรียนบทนี้จบแล้ว นักศึกษาจะสามารถ
\begin{enumerate}[topsep=0.2em, itemsep=0.15em, leftmargin=*]
    \item อธิบายโครงสร้างและประเภทของเมทริกซ์ได้
    \item ทำการบวก ลบ คูณเมทริกซ์ได้อย่างถูกต้อง
    \item คำนวณดีเทอร์มิแนนต์ของเมทริกซ์ขนาด 2×2 และ 3×3 ได้
    \item หาเมทริกซ์ผกผันได้
    \item ประยุกต์ใช้เมทริกซ์ในการแก้ระบบสมการเชิงเส้น การแปลงทางคอมพิวเตอร์กราฟิกส์ และการประมวลผลภาพได้
\end{enumerate}

\noindent \textbf{\large หัวข้อที่สอน}\par\vspace{0.2em}\noindent
\begin{enumerate}[topsep=0.2em, itemsep=0.15em, leftmargin=*]
    \item ความรู้เบื้องต้นเกี่ยวกับเมทริกซ์
    \item การดำเนินการบนเมทริกซ์
    \item ดีเทอร์มิแนนต์
    \item เมทริกซ์ผกผัน
    \item การประยุกต์ใช้เมทริกซ์ในคอมพิวเตอร์กราฟิกส์และการประมวลผลภาพ
\end{enumerate}

\noindent \textbf{\large สื่อการสอน}
\begin{enumerate}[topsep=0.2em, itemsep=0.15em, leftmargin=*]
    \item สไลด์ประกอบการสอน
    \item ตัวอย่างการคำนวณ
    \item แบบฝึกหัดท้ายบท
\end{enumerate}
